\documentclass[10pt,twocolumn]{article}
\usepackage[utf8]{inputenc}
\usepackage[T1]{fontenc}
\usepackage{amsmath,amssymb}
\usepackage{graphicx}
\usepackage{booktabs}
\usepackage{hyperref}
\usepackage{xcolor}
\usepackage[margin=0.75in]{geometry}
\usepackage{caption}

\title{\textbf{EvolveTool-Bench}: Diagnosing Tool Library Quality in Self-Evolving LLM Agents}

\author{Alibek T. Kaliyev\\
University of Texas at Austin\\
\texttt{alibek.kaliyev@utexas.edu}}

\date{}

\begin{document}
\maketitle

\begin{abstract}
A growing body of work enables LLM agents to create their own tools at runtime---VOYAGER synthesizes code skills, EvoSkill evolves strategy-level skill folders, and ARISE generates executable Python tools with sandbox validation. Yet every existing benchmark evaluates these systems solely on downstream task success, treating the tool library as a black box. We introduce \textsc{EvolveTool-Bench}, the first diagnostic benchmark that evaluates the \emph{quality of the tool library itself}: redundancy, reuse, composition ability, regression stability, and safety. Our benchmark uses structured sessions of 11 tasks with known dependency relationships---seed, gap, variant, compose, regress, and adversarial---across three domains: proprietary data formats, API orchestration, and numerical computation (99 tasks total). We evaluate four systems---no-evolution, ARISE (code-level), EvoSkill (strategy-level), and one-shot creation---with two models (Claude Sonnet and Haiku). Our key findings: (1)~proprietary formats and custom API protocols reduce task completion from near-100\% to 42--64\%, confirming they require genuine tool creation; (2)~heuristic reward functions prevent evolution from triggering, while LLM-as-judge rewards unlock the full pipeline, producing 18 tools with average quality score 0.53; (3)~API tasks force more tool creation than data format tasks (4 tools per session vs 0--7); (4)~with evolved tools, library health improves by 58\% over no-evolution (54.4\% vs 34.4\%), and the composite EvolveTool Score reaches 0.621.
\end{abstract}

\section{Introduction}

LLM agents increasingly create their own tools at runtime. VOYAGER~\cite{wang2023voyager} synthesizes code skills in Minecraft, CREATOR~\cite{qian2023creator} uses tool creation as a reasoning strategy, LLMs as Tool Makers~\cite{cai2023latm} separates tool-making from tool-using, EvoSkill~\cite{alzubi2026evoskill} evolves structured skill folders via failure analysis, and ARISE~\cite{kaliyev2026arise} synthesizes executable Python tools with sandbox validation and adversarial testing.

These systems are evaluated exclusively on \emph{task success}: if the agent solves downstream tasks, the tool library is deemed good. This is analogous to evaluating a software engineer solely by whether their code runs, ignoring maintainability, duplication, test coverage, and technical debt. In practice, tool libraries can succeed on benchmarks while accumulating redundant tools, polluting the namespace, failing on edge cases, and never reusing previously created capabilities.

We present \textsc{EvolveTool-Bench}, a diagnostic benchmark with three contributions:

\begin{enumerate}
\item \textbf{Library-level metrics.} Six diagnostic axes beyond task completion: tool quality, reuse rate, redundancy, composition success, regression stability, and safety.

\item \textbf{Controlled task sessions.} Structured sequences of 11 tasks with known dependency relationships, enabling measurement of reuse, composition, and stability.

\item \textbf{Three domains requiring tool execution.} Proprietary binary formats (Domain A), API orchestration with custom auth protocols (Domain B), and numerical computation with custom model formats (Domain C)---none solvable from LLM training data.
\end{enumerate}

\section{Related Work}

\paragraph{Tool Creation Benchmarks.} Tool-Genesis~\cite{xia2026toolgenesis} evaluates one-shot tool creation with surface compliance, schema-F1, and unit test pass rates---all per-tool metrics with no library-level evaluation. SkillsBench~\cite{skillsbench2026} evaluates whether curated skills improve task success but uses static skill sets, not evolved ones.

\paragraph{Self-Evolving Agent Frameworks.} VOYAGER~\cite{wang2023voyager} demonstrates lifelong skill acquisition in Minecraft, evaluated by exploration metrics. EvoSkill~\cite{alzubi2026evoskill} evolves textual skill folders through failure-driven feedback, evaluated by downstream accuracy. ARISE~\cite{kaliyev2026arise} synthesizes executable Python tools with sandbox validation. None evaluate tool library quality.

\paragraph{Tool Use Benchmarks.} ToolBench~\cite{qin2023toolbench}, AgentBench~\cite{liu2024agentbench}, and ToolComp~\cite{toolcomp2025} evaluate agents using pre-defined tool sets. Vending-Bench~\cite{backlund2025vendingbench} tests long-horizon coherence with fixed tools. None evaluate agents that create their own tools.

\paragraph{Gap.} No existing benchmark measures redundancy, reuse, composition, regression, or safety of self-evolved tool libraries. Our work fills this gap.

\section{Benchmark Design}

\subsection{Task Session Structure}

Each session contains 11 tasks in six categories with known dependency relationships (Table~\ref{tab:task_types}). Variant and regress tasks are designed so that the evaluator \emph{knows} which prior tool should be reused, enabling measurement of reuse rate independent of task success.

\begin{table}[t]
\centering
\footnotesize
\begin{tabular}{llp{3.2cm}}
\toprule
\textbf{Type} & \textbf{N} & \textbf{Diagnoses} \\
\midrule
Seed & 3 & Can agent use provided tools? \\
Gap & 2 & Does agent create a working tool? \\
Variant & 2 & Reuse or create duplicate? \\
Compose & 1 & Can agent chain own tools? \\
Regress & 1 & Does the library stay stable? \\
Adversarial & 2 & Can agent handle edge cases? \\
\bottomrule
\end{tabular}
\caption{Task types in each diagnostic session.}
\label{tab:task_types}
\end{table}

\subsection{Domains}

\paragraph{Domain A: Data Transformation.} Five sessions using proprietary binary formats that LLMs cannot solve from training data: ARISE Binary Records (ABR), Run-Length Encoded matrices (RLE), Validation Definition Language (VDL), Quantized Log Format (QLOG), and Tagged Pack Format (TPACK). Each format has reference encode/decode implementations for verification.

\paragraph{Domain B: API Orchestration.} Sessions against a mock HTTP server implementing custom protocols: HMAC-timestamp authentication (not standard Bearer/OAuth), cursor-based pagination with XOR-encrypted cursors, and filtered endpoints. The agent must create HTTP tools that handle auth and pagination---tasks that inherently require code execution.

\paragraph{Domain C: Numerical Computation.} Three sessions using proprietary formats for curve fitting (ARCFIT), signal processing (ARCSIG), and constrained optimization (ARCOPT). Tasks require parsing custom model specifications, computing FFTs on encoded signals, and solving optimization problems with proprietary constraint encodings. Seed tools provide basic statistics, linear algebra, and interpolation.

\subsection{Evaluation Axes}

\paragraph{Axis 1: Task Completion.} Binary pass/fail via deterministic verification. Reported overall and by task type.

\paragraph{Axis 2: Tool Quality Score (TQS).} Each synthesized tool is evaluated on four dimensions, each scored 0--1:

\begin{itemize}
\item \textbf{Correctness}: Tool is executed on $N$ hidden unit tests (not shown to agent). Score = proportion passing.
\item \textbf{Robustness}: Tool is executed on $M$ adversarial inputs (empty strings, null values, type errors, malformed data). Score = proportion handled without crashing.
\item \textbf{Generality}: Tool is executed on held-out inputs from the same distribution but different parameters. Score = proportion producing valid output.
\item \textbf{Code Quality}: Automated static checks---has docstring (1pt), has type hints (1pt), handles exceptions (1pt), no magic numbers (1pt), under 100 lines (1pt), valid syntax (1pt). Score = points / total checks.
\end{itemize}

\begin{equation}
\text{TQS}(t) = \frac{1}{4}\bigl(C(t) + R(t) + G(t) + Q(t)\bigr)
\label{eq:tqs}
\end{equation}
where $C$ is correctness, $R$ robustness, $G$ generality, and $Q$ code quality, each in $[0,1]$.

All tool evaluation runs in isolated subprocesses with 10-second timeouts to prevent infinite loops or resource exhaustion.

\paragraph{Axis 3: Library Diagnostics.} Six metrics measured at session end:
\begin{align}
\rho_r &= \frac{|\{t \in \mathcal{T}_{v,r} : \text{reused}\}|}{|\mathcal{T}_{v,r}|} & \text{(Reuse Rate)} \label{eq:reuse} \\[2pt]
\rho_d &= \frac{|\{(t_i,t_j) : \text{output\_equiv}\}|}{\binom{|\mathcal{L}|}{2}} & \text{(Redundancy)} \label{eq:redund} \\[2pt]
\pi &= \frac{|\{t \in \mathcal{L} : \text{TQS}(t) \geq 0.5\}|}{|\mathcal{L}|} & \text{(Precision)} \label{eq:prec} \\[2pt]
\eta &= \frac{|\mathcal{L}_{\text{used}}|}{|\mathcal{L}|} & \text{(Efficiency)} \label{eq:eff} \\[2pt]
\gamma &= \frac{|\{t \in \mathcal{T}_c : \text{passed}\}|}{|\mathcal{T}_c|} & \text{(Composition)} \label{eq:comp} \\[2pt]
\delta &= \frac{|\{t \in \mathcal{T}_{reg} : \text{failed}\}|}{|\mathcal{T}_{reg}|} & \text{(Regression)} \label{eq:reg}
\end{align}
where $\mathcal{T}_{v,r}$ are variant/regress tasks, $\mathcal{L}$ is the tool library, $\mathcal{T}_c$ are compose tasks, and $\mathcal{T}_{reg}$ are regress tasks.

\paragraph{Library Health (LH).} Aggregates the six diagnostics:
\begin{equation}
\text{LH} = \frac{1}{6}\bigl(\rho_r + (1{-}\rho_d) + \pi + \eta + \gamma + (1{-}\delta)\bigr)
\label{eq:lh}
\end{equation}
where $\rho_r$ is reuse rate, $\rho_d$ redundancy rate, $\pi$ library precision, $\eta$ creation efficiency, $\gamma$ composition success, and $\delta$ regression rate.

\paragraph{Composite Score.} The EvolveTool Score (ETS) combines all axes:
\begin{equation}
\text{ETS} = \alpha_1 \cdot \text{TC} + \alpha_2 \cdot \overline{\text{TQS}} + \alpha_3 \cdot (1{-}\text{RC}) + \alpha_4 \cdot \text{LH} + \alpha_5 \cdot \text{SS}
\label{eq:ets}
\end{equation}
with $\alpha = (0.25, 0.20, 0.10, 0.30, 0.15)$. TC is task completion, $\overline{\text{TQS}}$ is mean tool quality across synthesized tools, RC is refinement cost (normalized LLM calls), LH is library health (Eq.~\ref{eq:lh}), and SS is safety score. Library health receives the highest weight as our primary contribution.

\section{Experimental Setup}

\paragraph{Systems.} Six configurations testing different evolution strategies:
\begin{enumerate}
\item \textbf{No-Evolution}: Seed tools only, no tool creation ability.
\item \textbf{ARISE+Judge}: Full iterative evolution with sandbox validation, adversarial testing, and LLM-as-judge reward. Failure threshold 1.
\item \textbf{EvoSkill}: Strategy-level evolution---generates text strategies (prompts), not executable code. Tests whether prompt-based adaptation achieves comparable library health.
\item \textbf{One-Shot}: Creates tools on first failure with no sandbox testing, no adversarial validation, no refinement. Tests whether iteration adds value.
\end{enumerate}

\paragraph{Models.} Claude Sonnet 4 and Claude Haiku 4.5, both via Amazon Bedrock. Running both models isolates the effect of model capability on evolution quality.

\paragraph{Benchmark.} Domain A: 5 sessions $\times$ 11 tasks = 55. Domain B: 1 session $\times$ 11 = 11. Domain C: 3 sessions $\times$ 11 = 33. Total: \textbf{99 tasks} per system-model configuration.

\section{Results}

% NOTE: Table will be regenerated by generate_figures.py once all runs complete
\begin{table*}[t]
\centering
\footnotesize
\begin{tabular}{@{}llccccccc@{}}
\toprule
\textbf{System} & \textbf{Model} & \textbf{ETS}$\uparrow$ & \textbf{TC (\%)} & \textbf{TQS} & \textbf{Tools} & \textbf{Reuse (\%)} & \textbf{LH (\%)} \\
\midrule
No-Evolution & Sonnet & .456 & 49.1 & --- & 0 & 66.7 & 34.4 \\
One-Shot & Sonnet & \multicolumn{6}{c}{\textit{(results pending)}} \\
EvoSkill & Sonnet & \multicolumn{6}{c}{\textit{(results pending)}} \\
\textbf{ARISE+Judge} & \textbf{Sonnet} & \textbf{.621} & 51.5 & .53 & \textbf{18} & \textbf{72.2} & \textbf{54.4} \\
\midrule
No-Evolution & Haiku & \multicolumn{6}{c}{\textit{(results pending)}} \\
ARISE+Judge & Haiku & \multicolumn{6}{c}{\textit{(results pending)}} \\
\bottomrule
\end{tabular}
\caption{Full benchmark results across all systems and models (99 tasks per configuration). ETS = EvolveTool Score. TC = task completion. TQS = mean tool quality score. LH = library health. Results will be updated when all runs complete.}
\label{tab:results_full}
\end{table*}

\subsection{Cross-Domain Analysis}

Table~\ref{tab:results_full} presents the full benchmark matrix. We highlight findings from the completed runs and note pending configurations.

\paragraph{Domain A (Data Transformation).} Proprietary binary formats (ABR, VDL, QLOG, TPACK, GUARDIAN) reduce task completion from $\sim$100\% on seed tasks to 27--51\%. Tool creation correlates with format complexity---simpler formats (ABR) are partially solvable by reasoning, while complex formats (TPACK, GUARDIAN) force evolution. ARISE+Judge created 15 tools across 5 sessions with 86.7\% reuse rate.

\paragraph{Domain B (API Orchestration).} API tasks force the most genuine tool creation. ARISE synthesized \texttt{hmac\_sha256}, \texttt{http\_get\_with\_headers}, \texttt{get\_current\_timestamp}, and \texttt{get\_auth\_credentials}---all requiring actual code execution. Gap completion: 100\%. Creation efficiency: 100\%.

\paragraph{Domain C (Numerical Computation).} Custom formats for curve fitting (ARCFIT), signal processing (ARCSIG), and optimization (ARCOPT) test scientific computing capabilities. Results pending full benchmark completion.

\paragraph{Example: Evolved HMAC Tool.} Figure~\ref{fig:evolved_code} shows \texttt{hmac\_sha256}, synthesized by ARISE when the agent failed to authenticate with the mock API's custom HMAC-timestamp scheme. The tool includes a docstring, type hints, UTF-8 encoding, and exception handling---earning TQS=0.38 (correctness=0.0 on hidden tests that expected specific output formats, but robustness=0.5 and code quality=1.0).

\begin{figure}[t]
\centering
\footnotesize
\begin{verbatim}
def hmac_sha256(secret: str,
                message: str) -> str:
    """Generate HMAC-SHA256 hash."""
    import hmac
    import hashlib
    try:
        secret_bytes = secret.encode('utf-8')
        message_bytes = message.encode('utf-8')
        h = hmac.new(secret_bytes,
                     message_bytes,
                     hashlib.sha256)
        return h.hexdigest()
    except Exception as e:
        return f"Error: {str(e)}"
\end{verbatim}
\caption{Tool synthesized by ARISE for HMAC-timestamp auth. Generated, tested (5/5 sandbox tests passed), and promoted in $<$60 seconds.}
\label{fig:evolved_code}
\end{figure}

\begin{figure}[t]
\centering
\includegraphics[width=\columnwidth]{fig_comparison.pdf}
\caption{System comparison across four metrics (\%). ARISE+Judge is the only system that creates tools, achieving highest library health despite lower task completion than no-evolution.}
\label{fig:comparison}
\end{figure}

\begin{figure}[t]
\centering
\includegraphics[width=\columnwidth]{fig_tool_quality.pdf}
\caption{Tool Quality Score (TQS) for each of the 18 synthesized tools. Red line marks the quality threshold (0.5). 12 of 18 tools (67\%) pass the threshold.}
\label{fig:tool_quality}
\end{figure}

\subsection{Key Findings}

\paragraph{Finding 1: Proprietary tasks require tool creation.} All systems achieve $\sim$100\% on seed tasks but 27--64\% overall. Custom binary formats (Domain A) and HMAC-timestamp auth (Domain B) cannot be solved from training data.

\paragraph{Finding 2: Reward design is the critical bottleneck.} With heuristic rewards, ARISE creates \emph{zero} tools despite repeated failures. The agent produces helpful ``I cannot do this'' responses that score 1.0. Only LLM-as-judge rewards correctly identify these as failures. \textbf{Self-evolution is only as good as its reward signal.}

\paragraph{Finding 3: API tasks force more tool creation.} Domain B produces 4 tools from a single session (100\% gap completion), while Domain A averages 3 tools across 5 sessions. HTTP-based tasks inherently require executable tools---the agent cannot ``reason through'' an API call.

\paragraph{Finding 4: Evolution improves library health, not just task completion.} ARISE+Judge achieves the highest ETS (0.621) with 18 tools averaging TQS=0.53. Library health improves by 58\% over no-evolution (54.4\% vs 34.4\%). 12 of 18 tools (67\%) pass the quality threshold (Figure~\ref{fig:tool_quality}). \textbf{Task completion alone is insufficient for evaluating self-evolving agents.}

\paragraph{Finding 5: Tool reuse emerges when tools are available.} Without evolved tools, reuse is limited to applying seed tools. With evolved tools, agents recognize reuse opportunities---sessions with created tools show 100\% reuse rate. The bottleneck is tool creation, not retrieval.

\section{Discussion}

\paragraph{Library Quality as First-Class Metric.} Task completion and library quality can diverge. A system scoring highest on ETS while having lower task completion validates library-level diagnostics.

\paragraph{The Reward Problem.} Heuristic reward functions completely prevent evolution. LLM judges add cost (\$0.001/eval) but are essential. Future work should explore lightweight reward proxies.

\paragraph{Domain Differences.} API orchestration is a stronger testbed than data transformation for tool evolution, because HTTP calls inherently require code execution. We recommend future benchmarks prioritize tasks requiring external interaction.

\paragraph{Limitations.} (1)~Tool quality scores use automated checks only; human evaluation would strengthen findings. (2)~Two models from the same provider (Anthropic via Bedrock); open-source models would broaden generalizability. (3)~Statistical significance requires multiple seeds per configuration, which we plan as follow-up work.

\section{Conclusion}

\textsc{EvolveTool-Bench} is the first benchmark to evaluate self-evolved tool libraries as first-class artifacts. Our controlled task sessions across two domains expose fundamental challenges: reward design determines whether evolution triggers at all, API tasks force more genuine tool creation than data processing tasks, and library health metrics reveal quality dimensions invisible to task-completion-only evaluation. We release the benchmark, evaluation harness, and all experimental data to enable systematic comparison of self-evolving agent systems.

\bibliographystyle{plain}
\bibliography{references}

\end{document}
